% Generated draft: requires mathematical and source review.
\documentclass[11pt]{article}
\usepackage[margin=25mm]{geometry}
\usepackage{fontspec}
\setmainfont{Latin Modern Roman}
\usepackage{amsmath,amssymb,mathtools,unicode-math}
\setmathfont{Latin Modern Math}
\usepackage{enumitem,needspace,xurl,hyperref}
\setlength{\parindent}{0pt}
\setlength{\parskip}{5pt}
\setlength{\emergencystretch}{3em}
\newcommand{\N}{\mathbb{N}}
\newcommand{\Z}{\mathbb{Z}}
\newcommand{\Q}{\mathbb{Q}}
\newcommand{\R}{\mathbb{R}}
\newcommand{\C}{\mathbb{C}}
\newcommand{\F}{\mathbb{F}}
\newcommand{\A}{\mathbb{A}}
\newcommand{\PP}{\mathbb{P}}
\newcommand{\E}{\mathbb{E}}
\newcommand{\eps}{\varepsilon}
\newcommand{\dd}{\,\mathrm{d}}
\begin{document}
\section*{UnsolvedMath ID: 700849}
\textbf{Source problem code:} catalogue entry
\subsection{Definitions and mathematical statement}
For groups \(G\) and \(H\), write \(G*H\) for their free product and \(G\times H\) for their direct product. The free product is characterized by the universal property that every pair of homomorphisms from \(G\) and \(H\) into a group \(K\) extends uniquely to a homomorphism \(G*H\to K\). The direct product consists of pairs with componentwise multiplication.

\textit{Source question (O. N. Golovin; proposed in the supplied record by S. V. Ivanov).} On the class of all groups, does there exist an associative product operation satisfying the postulates of MacLane and Mal'cev, described in the source as being free, functorial, and hereditary, which is different from both the free-product operation and the direct-product operation?

\textit{Unresolved qualifications.} The supplied record does not state these postulates or define the three properties in their intended technical sense. Consequently, they remain requirements referring to the source terminology, rather than explicit axioms in this draft. The record also does not specify whether associativity means equality, isomorphism, or a coherent natural isomorphism, nor which equivalence of product operations is intended by ``different.'' Thus the source question is preserved, but a fully specified mathematical existence statement requires the original postulates and conventions.
\subsection{Short English statement}
The question asks whether the class of all groups admits an associative product satisfying the specified MacLane--Mal'cev postulates other than free and direct products. The supplied record omits the postulates and the conventions for comparing operations, so these require clarification before a precise axiomatic formulation is possible.
\subsection{Sources}
\begin{itemize}
\item Local UnsolvedMath problem registry, permanent ID \(700849\), exact source code \texttt{catalogue entry}, title \emph{Kourovka Notebook Problem 12.30}. The supplied registry labels the problem \emph{open}; this label has not been independently verified.
\item \emph{The Kourovka Notebook (Unsolved Problems in Group Theory)}, 12th issue (1992), Problem 12.30; also identified in the supplied background with the 21st issue edition (2026), \texttt{arXiv:1401.0300v48}. The collection description names E. I. Khukhro and V. D. Mazurov as editors. This bibliographic identification is supplied evidence; the referenced text has not been consulted.
\item Attribution: the statement begins with O. N. Golovin's name, while the background and proposer field name S. V. Ivanov; the supplied proposal year is 1992. Both attributions are retained without inferring their precise relationship.
\end{itemize}
\end{document}
